%_ **************************************************************************
%_ * The TeX source for AMS journal articles is the publishers TeX code     *
%_ * which may contain special commands defined for the AMS production      *
%_ * environment.  Therefore, it may not be possible to process these files *
%_ * through TeX without errors.  To display a typeset version of a journal *
%_ * article easily, we suggest that you retrieve the article in DVI,       *
%_ * PostScript, or PDF format.                                             *
%_ **************************************************************************
%  Author Package
%% Translation via Omnimark script a2l, October 5, 1999 (all in one day!)
%% Declarations:
%% User definitions:


\documentclass{proc-l}
\usepackage{amssymb}

%%%%%%%%%%%%%%%%%%%%%%%%%%%%%%%%%%%%%%%%%%%%%%%%%%%%%%%%%%%%%%%%%%%%%%%%%%%%%%%%%%%%%%%%%%%%%%%%%%%%
%TCIDATA{OutputFilter=LATEX.DLL}
%TCIDATA{LastRevised=Thursday, July 12, 2001 07:52:20}
%TCIDATA{<META NAME="GraphicsSave" CONTENT="32">}

\controldates{20-OCT-2000,20-OCT-2000,20-OCT-2000,23-OCT-2000}
\issueinfo{129}{06}{June}{2001}
\pagespan{1727}{1731}
\dateposted{October 31, 2000}
\PII{S 0002-9939(00)05701-4}
\copyrightinfo{2000}{American Mathematical Society}
\theoremstyle{plain}
\newtheorem*{theorem1}{Theorem 1}
\newtheorem*{theorem2}{Theorem 3}
\newtheorem*{theorem3}{Theorem 5}
\theoremstyle{remark}
\newtheorem*{remark1}{Remark 2}
\theoremstyle{definition}
\newtheorem*{definition1}{Example 4}
\newcommand{\CC}{\mathbb{C}}
\newcommand{\myL}{\mathcal{L}}
\newcommand{\R}{\mathcal{R}}
\newcommand{\N}{\mathcal{N}}
\newcommand{\ind}{\operatorname {ind}}
\newcommand{\dplus}{\overset \bullet \to {+}}
\newcommand{\ov}[1]{\overline {#1}}

\input{tcilatex}

\begin{document}
\title{Products of EP operators on Hilbert spaces}
\author{Dragan S. Djordjevi\'{c}}
\address{Department of Mathematics, Faculty of Philosophy, University of Ni%
\v{s}, \'{C}irila i\break Metodija 2, 18000 Ni\v{s}, Yugoslavia}
\email{dragan@archimed.filfak.ni.ac.yu}
\email{dragan@filfak.filfak.ni.ac.yu}
\date{May 4, 1999 and, in revised form, September 17, 1999}
\subjclass[2000]{Primary 47A05, 15A09}
\keywords{EP operators, generalized inverses}

\begin{abstract}
A Hilbert space operator $A$ is called the EP operator if the range of $A$
is equal to the range of its adjoint $A^{*}$. In this article necessary and
sufficient conditions are given for a product of two EP operators with
closed ranges to be an EP operator with a closed range. Thus, a
generalization of the well-known result of Hartwig and Katz (Linear Algebra
Appl. \textbf{252 } (1997), 339--345) is given.
\end{abstract}

\maketitle

\commby{Joseph A. Ball}

\section*{1. Introduction and preliminaries}

Let $H,H_{1},H_{2}$ be arbitrary Hilbert spaces, $\mathcal{L} (H_{1},H_{2})$
the set of all  bounded operators from $H_{1}$ into $H_{2}$ and $\mathcal{L}
(H,H)=\mathcal{L} (H)$.  For $A\in \mathcal{L} (H_{1},H_{2})$ we use $A^{*}$%
, $\mathcal{R} (A)$ and $\mathcal{N} (A)$, respectively, to denote the
adjoint operator, range and kernel of $A$.

If $\mathcal{R} (A)$ is closed, then $A^{\dag }\in \mathcal{L} (H_{2},H_{1})$
is defined as the unique operator in $\mathcal{L} (H_{2},H_{1})$ which
satisfies the equations: 
\begin{equation*}
AA^{\dag }A=A,\quad A^{\dag }AA^{\dag }=A^{\dag },\quad (AA^{\dag
})^{*}=AA^{\dag },\quad (A^{\dag }A)^{*}=A^{\dag }A.
\end{equation*}
$A^{\dag }$ is known as the Moore-Penrose inverse of $A$ \cite{2}, \cite{4}.

For $A\in \mathcal{L} (H)$ consider the following inclusions: $\{0\}\subset%
\mathcal{N} (A)\subset\mathcal{N} (A^{2})\subset\cdots $ and $H\supset 
\mathcal{R} (A)\supset \mathcal{R} (A^{2})\supset \cdots $. The ascent of $A$
is defined as the least $k$ (if it exists) for which the following holds: $%
\mathcal{N} (A^{k})=\mathcal{N} (A^{k+1})$. If such $k$ does not exist, we
say that the ascent of $A$ is equal to infinity. The descent of $A$ is
defined as the least $k$ (if it exists) for which the following is
satisfied: $\mathcal{R} (A^{k})=\mathcal{R} (A^{k+1})$. If such $k$ does not
exist, we say that the descent of $A$ is equal to infinity. If the ascent
and the descent of $A$ are finite, then they are equal, and this common
value is known as the Drazin index of $A$, denoted by $\func{ind} (A)$. If $%
\func{ind} (A)=k$, then there exists the unique operator $A^{D}\in \mathcal{L%
} (H)$, which satisfies 
\begin{equation*}
A^{k+1}A^{D}=A^{k},\quad A^{D} AA^{D}=A^{D},\quad AA^{D}=A^{D}A.
\end{equation*}
$A^{D}$ is known as the Drazin inverse of $A$. If $\func{ind} (A)=0$, then $A
$ is invertible. If $\func{ind} (A)\leq 1$, then $A^{D}$ is known as the
group inverse of $A$, denoted by $A^{\#}$ \cite{2}, \cite{4}.

In this article we shall consider operators which satisfy $A^{D}=A^{\dag }$.
Notice that the last condition implies $\func{ind} (A)\leq 1$; hence we
shall consider operators which satisfy $A^{\#}=A^{\dag }$.

Recall that an operator $A\in \mathcal{L} (H)$ is called the EP operator if $%
\mathcal{R} (A)=\mathcal{R} (A^{*})$ \cite{1}, \cite{5}, \cite{6}.

If $A\in \mathcal{L} (H)$ has a closed range, then it is easy to see that $A$
is an EP operator if and only if $A^{\dag }=A^{\#}$. Also, if $A$ is an EP
operator, then $\mathcal{N} (A)=\mathcal{N} (A^{*})$.

If $A$ and $B$ are two EP operators, it is a question when $AB$ is the EP
operator. This problem for complex square matrices was open for twenty-five
years (see \cite{1}, \cite{5}). In \cite{5} Hartwig and Katz found necessary
and sufficient conditions for a product of two square EP matrices to also be
the EP matrix. Their method cannot be applied to arbitrary Hilbert space
operators.

In this paper we give necessary and sufficient conditions for a product of
two EP operators with closed ranges to be the EP operator with a closed
range. Thus, using a different method, we generalize the result from \cite{5}%
.

We shall mention the matrix decompositions of $A^{\dag }$ and $A^{\#}$. If $M
$ and $N$ are subspaces of $H$, then $M\overset{\bullet}{+} N$ denotes the
direct (not necessary orthogonal) sum of $M$ and $N$. However, $M\oplus N$
denotes the orthogonal sum of $M$ and $N$.

If $A\in \mathcal{L} (H)$ has the group inverse, then $H=\mathcal{R} (A)%
\overset{\bullet}{+} \mathcal{N} (A)$ and $A$ has the following matrix
decomposition: 
\begin{equation*}
A=%
\begin{bmatrix}
A_{1} & 0 \\ 
0 & 0%
\end{bmatrix}
:%
\begin{bmatrix}
\mathcal{R} (A) \\ 
\mathcal{N} (A)%
\end{bmatrix}
\to 
\begin{bmatrix}
\mathcal{R} (A) \\ 
\mathcal{N} (A)%
\end{bmatrix}
,  \tag{1}
\end{equation*}
where $A_{1}=A|_{\mathcal{R} (A)}:\mathcal{R} (A)\to \mathcal{R} (A)$ is
invertible. In this case the group inverse of $A$ has the form 
\begin{equation*}
A^{\#}=%
\begin{bmatrix}
A_{1}^{-1} & 0 \\ 
0 & 0%
\end{bmatrix}
.
\end{equation*}

If $A\in \mathcal{L} (H_{1},H_{2})$ and $\mathcal{R} (A)$ is closed, then $A$
has the following matrix decomposition with  respect to the orthogonal sums
of subspaces of $H_{1}$ and $H_{2}$: 
\begin{equation*}
A=\left[%
\begin{matrix}
{}\,\overline {A} & 0 \\ 
0 & 0%
\end{matrix}%
\right] :H_{1}=%
\begin{bmatrix}
\mathcal{R} (A^{*}) \\ 
\mathcal{N} (A)%
\end{bmatrix}
\to 
\begin{bmatrix}
\mathcal{R} (A) \\ 
\mathcal{N} (A^{*})%
\end{bmatrix}
=H_{2},
\end{equation*}
where $\overline {A}=A|_{\mathcal{R} (A^{*})}:\mathcal{R} (A^{*})\to 
\mathcal{R} (A)$ is invertible. The Moore-Penrose inverse of $A$ is given as 
\begin{equation*}
A^{\dag }=\left[%
\begin{matrix}
{}\,\overline {A}\,^{-1} & 0 \\ 
0 & 0%
\end{matrix}%
\right] :%
\begin{bmatrix}
\mathcal{R} (A) \\ 
\mathcal{N} (A^{*})%
\end{bmatrix}
\to 
\begin{bmatrix}
\mathcal{R} (A^{*}) \\ 
\mathcal{N} (A)%
\end{bmatrix}%
.
\end{equation*}

If $A\in \mathcal{L} (H)$ is an EP operator with a closed range, then (1)
becomes the decomposition of $A$ with respect to the orthogonal sum of
subspaces.

\section*{2. Results}

Now, we state the main result of this paper.

\begin{theorem1}
Let $A,B\in \mathcal{L} (H)$ be EP operators with closed ranges.  Then the
following statements are equivalent.

\begin{itemize}
\item[(a)] $AB$ is an EP operator with a closed range.

\item[(b)] $\mathcal{R} (AB)=\mathcal{R} (A)\cap \mathcal{R} (B)$ and $%
\mathcal{N} (AB)=\mathcal{N} (A)+\mathcal{N} (B)$.

\item[(c)] $\mathcal{R} (AB)=\mathcal{R} (A)\cap \mathcal{R} (B)$ and $%
\mathcal{N} (AB)=\overline{\mathcal{N} (A)+ \mathcal{N} (B)}$.
\end{itemize}
\end{theorem1}

\begin{proof} (a)$\implies $(b): Suppose that $AB$ is an EP operator 
with a
closed range.
Since $A$ is
an EP operator,  we get that the decomposition of $A$,
\begin{equation*}A=\begin{bmatrix}A_{1}&0\\0&0\end{bmatrix}
:\begin{bmatrix}\R (A)\\\N (A)\end{bmatrix}
\to \begin{bmatrix}\R (A)\\\N (A)\end{bmatrix},
\end{equation*}
is a decomposition with respect to the orthogonal sum of subspaces.
Since $B$ is an EP operator,  it follows that $B$ has the following
orthogonal decomposition:
\begin{equation*}B=\begin{bmatrix}B_{1}&0\\ B_{2}&0\end{bmatrix}
:\begin{bmatrix}\R (B)\\\N (B)\end{bmatrix}
\to \begin{bmatrix}\R (A)\\\N (A)\end{bmatrix}
.\end{equation*}
We can write $\R (B)=\N (B_{1})\oplus \N (B_{1})^{\bot }$.
It is easy to see $\N (AB)=\N (B)\oplus \N (B_{1})$, so
\begin{equation*}\R ((AB)^{*})=\N (AB)^{\bot }=\N (B_{1})^{\bot }\subset\R 
(B).\end{equation*}
Also, $\R (AB)=\R (A_{1}B_{1})=A_{1}(\R (B_{1}))$. Since $AB$ is an EP 
operator we
get $\N (B_{1})^{\bot }=\R (AB)$.

First, we prove $\R (AB)=\R (A)\cap \R (B)$. From
$\R (AB)=\N (B_{1})^{\bot }\subset\R (B)$ and $\R (AB)\subset\R (A)$, it 
follows that
$\N (B_{1})^{\bot }=\R (AB)\subset\R (A)\cap \R (B)\subset\R (B_{1})$.

Let $Q$ be a projection from $H$ onto $\R (A)$ parallel to $\N (A)$ and
let $P$ be a projection from $H$ onto $\R (B)$. By \cite[Theorem 1, 
p.~127]{4}, we know that $\R (QP)$ is closed if and only if
$\R (AB)$ is closed. Hence,  $\R (QP)=Q(\R (B))=\R (B_{1})$ is closed.

Suppose that there exists a closed subspace $S$ of $\R (B_{1})$, such 
that
\begin{equation*}\N (B_{1})^{\bot }\oplus S=\R (B_{1}).\end{equation*}
We conclude  $S\subset\N (B_{1})\oplus \N (B)$.

Define an operator $\ov {B_{1}}\in \myL (H)$ in the following way:
$\ov {B_{1}}|_{\R (B)}=B_{1}$ and $\ov {B_{1}}|_{\N (B)}\linebreak
=0$. Then
$\R (\ov {B_{1}})=\R (B_{1})$ and $\N (\ov {B_{1}})=\N (B)\oplus \N 
(B_{1})$.

Notice the following:
\begin{equation*}\R (\ov {B_{1}}^{2})=\ov {B_{1}}(\R (B_{1}))=B_{1}(\N 
(B_{1})^{\bot })=\R (B_{1})=
\R (\ov {B_{1}})\end{equation*}
and
\begin{equation*}\N (\ov {B_{1}}^{2})=[\N (B)\oplus \N (B_{1})]+S=\N 
(\ov {B_{1}}).\end{equation*}
It follows that the Drazin index of $\ov {B_{1}}$ is equal to
zero or one; hence
\begin{equation*}H=\N (\ov {B_{1}})\overset{\bullet}{+} \R (\ov {B_{1}})=[\N 
(B)\oplus \N (B_{1})]\overset{\bullet}{+} \R (B_{1}).\end{equation*}
On the other hand, since $S\subset[\N (B)\oplus \N (B_{1})]\cap \R (B_{1})$
we conclude
$S=\{0\}$. Finally, $\N (B_{1})^{\bot }=\R (A)\cap \R (B)=\R (B_{1})$ and
$\R (AB)=\R (A)\cap \R (B)$.

To prove $\N (AB)=\N (A)+\N (B)$, consider the following decompositions:
\begin{equation*}H=\N (B_{1})^{\bot }\oplus \N (B_{1})\oplus \N (B)=\R 
(B_{1})\oplus \R (B_{1})^{\bot }\oplus \N (A).\tag{2}\end{equation*}
We conclude
\begin{equation*}\N (AB)=\N (B)\oplus \N (B_{1})=\N (A)\oplus \R 
(B_{1})^{\bot }.\end{equation*}
From $\R (B)\subset\R (B_{1})\oplus \R (B_{2})$ we get
that $\R (B_{1})^{\bot 
}\subset\N (B)$, and
 it follows  $\N (AB)=\N (A)+\N (B)$.

(b)$\implies $(c): This part of the proof is obvious.

(c)$\implies $(a): Suppose that $\R (AB)=\R (A)\cap \R (B)$ and
$\N (AB)=\overline{\N (A)+\N (B)}$. Immediately we get that $\R (AB)$ is 
closed.
If $M$ and $N$ are subspaces of $H$, then\linebreak
$(M\cap N)^{\bot 
}=\overline{M^{\bot }+N^{\bot }}$.
Now we find
\begin{equation*}\begin{split}
\R (AB)&=(\R (A)\cap \R (B)){^{\bot }}{^{\bot }}=(\overline{(\R 
(A)^{\bot }+\R (B)^{\bot })}^{\bot }\\
&= (\overline{\N (A)+\N (B)})^{\bot }=\N (AB)^{\bot }=\R (AB)^{*},
\end{split}\end{equation*}
and finally,  $AB$ is an EP operator with a closed range.\end{proof}

Our Theorem 1 is a generalization of the main result of Hartwig and Katz %
\cite{5}.

\begin{remark1}
The statements (b) and (c) of Theorem 1 illustrate the difference between
the finite and infinite dimensional Hilbert spaces. If $H$ is a finite
dimensional Hilbert space, then the statement (b) is the same as the
statement (c) of Theorem 1. In general Hilbert spaces, the statement (b) is
a stronger condition than the statement (c) of Theorem 1.
\end{remark1}

It is well-known that the product of two selfadjoint operators $A,B$ is a
selfadjoint operator if and only if $AB=BA$. In \cite{6} it is shown that
the product of two commuting EP matrices is again the EP matrix. For EP
operators on an arbitrary Hilbert space the following result holds.

\begin{theorem2}
If $A,B\in \mathcal{L} (H)$ are EP operators with closed ranges and $AB=BA$,
then $AB$ is the EP operator with a closed range also.
\end{theorem2}

\begin{proof} It is well-known that the group inverse $A^{\#}$ of $A$
 commutes with every operator that commutes with $A$ \cite{4}, so
we conclude that operators $A,B, A^{\#}, B^{\#}$ mutually commute.
Obviously, $(AB)^{\#}=B^{\#}A^{\#}=B^{\dag }A^{\dag }$. Notice also
\begin{equation*}\begin{split}
(ABB^{\#}A^{\#})^{*}&=(BB^{\#}AA^{\#})^{*}=(BB^{\dag }AA^{\dag 
})^{*}=(AA^{\dag })^{*}(BB^{\dag })^{*}=AA^{\dag }BB^{\dag }\\
&=ABB^{\#} A^{\#}.
\end{split}\end{equation*}
The equality $(B^{\#}A^{\#} AB)^{*}=B^{\#}A^{\#} AB$ can be proved 
analogously. We
have just proved
\begin{equation*}(AB)^{\#}=B^{\#}A^{\#}=B^{\dag }A^{\dag }=(AB)^{\dag 
};\end{equation*}
hence $AB$ is an EP operator with a closed range.\end{proof}

We can construct two noncommuting EP operators whose product is also an EP
operator. Let $P$ be any orthogonal projection and let $A$ be any invertible
operator which satisfies $A(\mathcal{N} (P))=\mathcal{N} (P)$ and $A(%
\mathcal{R} (P))\nsubseteq \mathcal{R} (P)$. Then $P=P^{\dag }=P^{\#}$ and $%
A^{-1}=A^{\dag }=A^{\#}$. Also, 
\begin{equation*}
\mathcal{R} (P)\cap \mathcal{R} (A)=\mathcal{R} (P)=\mathcal{R} (PA),\quad 
\mathcal{N} (P)+\mathcal{N} (A)=\mathcal{N} (P)=\mathcal{N} (PA).
\end{equation*}
By Theorem 1 it follows that $PA$ is an EP operator. On the other hand, $A$
and $P$ commute if and only if the decomposition $H=\mathcal{R} (P)\oplus 
\mathcal{N} (P)$ reduces $A$, i.e. $A(\mathcal{R} (P))\subset\mathcal{R} (P)$
and $A(\mathcal{N} (P))\subset\mathcal{N} (P)$. According to our
construction we get $AP\neq PA$.

For a concrete realization of the preceding situation, consider the
following example.

\begin{definition1}
Let us take $H=\mathbb{C} ^{2}$, $P=%
\begin{bmatrix}
1 & 0 \\ 
0 & 0%
\end{bmatrix}
$ and $A=%
\begin{bmatrix}
1 & 0 \\ 
1 & 1%
\end{bmatrix}
$. Then $P=P^{\dag }=P^{\#}$ and $A^{-1}=A^{\dag }=A^{\#}=%
\begin{bmatrix}
1 & 0 \\ 
-1 & 1%
\end{bmatrix}
$; hence $P$ and $A$ are EP operators. Notice 
\begin{equation*}
PA=P\neq 
\begin{bmatrix}
1 & 0 \\ 
1 & 0%
\end{bmatrix}
=AP.
\end{equation*}
Now, $P$, $A$ and $PA$ are EP operators, but $AP\neq PA$.
\end{definition1}

If $\mathcal{R} (A)$, $\mathcal{R} (B)$ and $\mathcal{R} (AB)$ are closed,
then the rule $(AB)^{\dag }=B^{\dag }A^{\dag }$ is called the reverse order
rule for the Moore-Penrose inverse (and it does not hold in general). In %
\cite{3} it is shown that if $\mathcal{R} (A)$, $\mathcal{R} (B)$ and $%
\mathcal{R} (AB)$ are closed, then the following statements are equivalent:

\begin{itemize}
\item[(a)] $(AB)^{\dag }=B^{\dag }A^{\dag }$;

\item[(b)] $\mathcal{R} (A^{*}AB)\subset\mathcal{R} (B)$ and $\mathcal{R}
(BB^{*}A^{*})\subset \mathcal{R} (A^{*})$.
\end{itemize}

In \cite{1} Baskett and Katz considered the connection between the EP
matrices and the reverse order rule for the Moore-Penrose inverse of
matrices. To illustrate this connection, we formulate the following result.

\begin{theorem3}
If $A,B\in \mathcal{L} (H)$ are EP operators with closed ranges and $%
\mathcal{R} (A)=\mathcal{R} (B)$, then $(AB)^{\dag }=B^{\dag }A^{\dag }$.
\end{theorem3}

\begin{proof} From $\R (A^{*})=\R (A)=\R (B)=\R (B^{*})$ we get
\begin{equation*}\R (A^{*}AB)=A^{*}A(\R (A^{*}))=A^{*}(\R (A))=\R 
(A^{*})=\R (B)\end{equation*}
and
\begin{equation*}\R (BB^{*}A^{*})=BB^{*}(\R (B))=\R (B)=\R 
(A^{*}).\end{equation*}
Hence, the reverse order rule $(AB)^{\dag }=B^{\dag }A^{\dag }$ 
holds.\end{proof}

Some other similar and interesting results can be found in \cite{1}.

\section*{Acknowledgement}

I am grateful to the referee for helpful comments and suggestions concerning
the paper.

\begin{thebibliography}{1}
\bibitem[1]{1} T.~S.~Baskett and I.~J.~Katz, \emph{Theorems on products of $%
EP_{r}$ matrices}, Linear Algebra Appl. \textbf{2} (1969), 87--103. %
\MR{40:4280}

\bibitem[2]{2} A.~Ben-Israel and T.~N.~E.~Greville, \emph{Generalized
inverses: theory and applications}, Wiley-In\-ter\-sci\-en\-ce, New York,
1974. \MR{53:469}

\bibitem[3]{3} R. H. Bouldin, \emph{Generalized inverses and factorizations}%
, Recent applications of generalized inverses, Pitman Ser. Res. Notes in
Math., vol.~66, 1982, pp.~233--249. \MR{83j:47001}

\bibitem[4]{4} S.~R.~Caradus, \emph{Generalized inverses and operator theory}%
, Queen's paper in pure and applied mathematics, Queen's University,
Kingston, Ontario, 1978. \MR{81m:47003}

\bibitem[5]{5} R.~E.~Hartwig and I.~J.~Katz, \emph{On products of EP matrices%
}, Linear Algebra Appl. \textbf{252} (1997), 339--345. \MR{98a:15050}

\bibitem[6]{6} I.~J.~Katz, \emph{Weigmann type theorems for $EP_{r}$ matrices%
}, Duke Math. J. \textbf{32} (1965), 423--427. \MR{31:4804}
\end{thebibliography}

\end{document}
